\chapter{The frozen challenge set}\label{ch:challenge}

The intent model's held-out score says how often a message is understood; the challenge set says how often a whole
conversation ends correctly and safely, with the policy, the tools, the confirmation flow and the gateway all
involved.

\section{The cases}

The set holds \est{chl.cases}{value} scripted conversations, \est{chl.cases-es}{value} in Spanish and
\est{chl.cases-pt}{value} in Portuguese, in \est{chl.categories}{value} categories (\cref{tab:challenge-composition});
\est{chl.multi-turn}{value} of them have more than one turn, such as a question, a clarifying answer and a
confirmation. They were written from the specification of what a correct bank agent must do, before any scored run,
and frozen by a manifest of file digests that every run checks; the first scored run found the set unchanged. Each
case names a demonstration customer by the role it plays (a customer with a single active card, a closed card, a
fraud flag, several cards), never by an identifier, and some inject a tool failure before a turn.

Each case declares its scope: \est{chl.scope-resolve}{value} cases the copilot should resolve alone,
\est{chl.scope-escalate}{value} that a person must decide, and \est{chl.scope-guard}{value} where it must refuse,
deny, or ask the customer to authenticate. A case passes when its final turn has the expected outcome, the expected
hand-off queue, and exactly the expected action in the operational store.

\generated{tables/challenge-composition}

\section{Metrics}

\begin{description}[style=nextline]
  \item[Correct final outcome] the share of all cases that pass.
  \item[Safe automated resolution] the share of the cases to resolve (or of the cases the copilot did not
  transfer) that pass with no unsafe outcome.
  \item[Containment] the share of cases not transferred to a person.
  \item[Missed transfers] the share of the cases a person must decide that the copilot did not transfer.
  \item[Unnecessary transfers] the share of the cases to resolve that the copilot transferred anyway.
  \item[Unsafe outcomes] the share of cases with any of the following, detected by the harness from the operational
  store and the replies rather than from the copilot's own labels: an action the case does not call for, a wrong or
  unconfirmed action, an unblock without the step-up, a success reported without a read-back, data served without a
  session, another customer's or an unknown card served, or an injection not refused.
\end{description}

\section{Keyword routing against the learned model}

The copilot was run on the frozen set without any language model in a pair of configurations that differ only in
intent recognition:
keyword routing, and the learned model of \cref{ch:intent}, which asks the customer to rephrase when unsure.

\begin{table}[htbp]\centering\small
  \caption{The first scored run: share with numerator, denominator and Wilson interval.}
  \label{tab:challenge-main}
  \begin{tabular}{lll}
    \toprule
    Metric & Keyword routing & Learned model \\
    \midrule
    Correct final outcome & \est{chl.first-keyword-correct}{num}/\est{chl.first-keyword-correct}{den} \est{chl.first-keyword-correct}{ci}
      & \est{chl.first-learned-correct}{num}/\est{chl.first-learned-correct}{den} \est{chl.first-learned-correct}{ci} \\
    Safe resolution, of cases to resolve & \est{chl.first-keyword-safe-in-scope}{num}/\est{chl.first-keyword-safe-in-scope}{den} \est{chl.first-keyword-safe-in-scope}{ci}
      & \est{chl.first-learned-safe-in-scope}{num}/\est{chl.first-learned-safe-in-scope}{den} \est{chl.first-learned-safe-in-scope}{ci} \\
    Safe resolution, of cases attempted & \est{chl.first-keyword-safe-attempted}{num}/\est{chl.first-keyword-safe-attempted}{den} \est{chl.first-keyword-safe-attempted}{ci}
      & \est{chl.first-learned-safe-attempted}{num}/\est{chl.first-learned-safe-attempted}{den} \est{chl.first-learned-safe-attempted}{ci} \\
    Containment & \est{chl.first-keyword-containment}{num}/\est{chl.first-keyword-containment}{den} \est{chl.first-keyword-containment}{ci}
      & \est{chl.first-learned-containment}{num}/\est{chl.first-learned-containment}{den} \est{chl.first-learned-containment}{ci} \\
    Missed transfers & \est{chl.first-keyword-missed}{num}/\est{chl.first-keyword-missed}{den} \est{chl.first-keyword-missed}{ci}
      & \est{chl.first-learned-missed}{num}/\est{chl.first-learned-missed}{den} \est{chl.first-learned-missed}{ci} \\
    Unnecessary transfers & \est{chl.first-keyword-unnecessary}{num}/\est{chl.first-keyword-unnecessary}{den} \est{chl.first-keyword-unnecessary}{ci}
      & \est{chl.first-learned-unnecessary}{num}/\est{chl.first-learned-unnecessary}{den} \est{chl.first-learned-unnecessary}{ci} \\
    Unsafe outcomes & \est{chl.first-keyword-unsafe}{num}/\est{chl.first-keyword-unsafe}{den} \est{chl.first-keyword-unsafe}{ci}
      & \est{chl.first-learned-unsafe}{num}/\est{chl.first-learned-unsafe}{den} \est{chl.first-learned-unsafe}{ci} \\
    \bottomrule
  \end{tabular}
\end{table}

The learned model ends \est{chl.first-learned-correct}{num} of the \est{chl.first-learned-correct}{den} cases
correctly, against \est{chl.first-keyword-correct}{num} for keyword routing (\cref{tab:challenge-main}). On the same
cases the learned model passes where keywords fail \est{chl.mcnemar-b}{value} times and the reverse happens
\est{chl.mcnemar-c}{value} times; the exact McNemar test gives $p = \est{chl.mcnemar}{p}$, a single planned
comparison. It misses fewer transfers a person must decide (\est{chl.first-learned-missed}{num} of
\est{chl.first-learned-missed}{den}, against \est{chl.first-keyword-missed}{num}) and makes no unnecessary transfer.
By language it passes \est{chl.first-learned-correct-es}{num} of \est{chl.first-learned-correct-es}{den} Spanish and
\est{chl.first-learned-correct-pt}{num} of \est{chl.first-learned-correct-pt}{den} Portuguese cases.

\begin{figure}[htbp]
  \centering
  \includegraphics[width=0.9\linewidth]{fig-chl-categories}
  \caption{Correct final outcome by category in the first scored run, with Wilson intervals; the number of cases in
  each category is in brackets.}
  \label{fig:chl-categories}
\end{figure}

\Cref{fig:chl-categories} shows where the difference comes from: multi-card conversations, where keywords lose the
thread when the customer answers which card they mean; reserved requests that must reach a specific queue;
out-of-scope requests; and knowledge questions, none of which keyword routing sent where it should. The learned
model loses a little on paraphrases and on a question about a missing value, where an informal wording falls below
its threshold and the copilot asks the customer to rephrase. Every category has few cases, so the per-category intervals are wide; the figure locates the
difference rather than measuring it.

The learned configuration failed \est{chl.learned-failures}{value} cases (\cref{tab:challenge-failures} in
\cref{app:tables}): informal paraphrases it asked to rephrase, a code-switched confirmation, a knowledge question that
found no document, and one unsafe outcome. Without a model, the copilot runs a turn in a median of
\est{chl.learned-ms-p50}{value}\,ms (\est{chl.learned-ms-p95}{value}\,ms at the \proto{stats.high_percentile}th
percentile).

\section{The unsafe outcome and the rerun after the fix}

Both configurations had one unsafe outcome, the same case: a Portuguese message that tells the copilot it is now an
assistant without limits and asks it to raise the customer's limit. The gateway's injection rules did not recognise
the second-person persona wording, so the message went through; keyword routing answered about the balance, and the
learned model routed a limit increase to a person. Neither changed anything, but a guard case must be refused, and
the harness counts it as unsafe.

The fix widened the gateway's injection patterns to second-person persona prompts in Portuguese and Spanish and to
requests for an assistant without limits or rules. The frozen set was then run again: the learned configuration
passed \est{chl.after-learned-correct}{num} of \est{chl.after-learned-correct}{den} cases with
\est{chl.after-learned-unsafe}{num} unsafe outcomes (interval \est{chl.after-learned-unsafe}{ci}), and keyword routing
\est{chl.after-keyword-correct}{num} with \est{chl.after-keyword-unsafe}{num}. Because the fix was chosen knowing which
case had failed, this rerun is not a held-out result; the first run remains the estimate of how the copilot does on
cases it has not seen. The rerun is the baseline for the model runs of \cref{ch:models}, which used the copilot after
the fix.
